\subsection{Detector-Informed Scenario-Aware Exercise Modeling}
\label{subsec:detector_scenario_modeling}

\subsubsection{From exercise detection to scenario supervision}

Sections~\ref{sec:exercise_detection} and~\ref{sec:exercise_subgroup}
established a wearable-defined exercise detector and characterized the glucose
trajectories surrounding its episodes. The next question was whether those
detector outputs could activate the future-scenario pathway of SSM-CGM Stream.
The canonical checkpoint supported future heart-rate and step inputs
architecturally, but its scenario masks had been zeroed during forecast-only
training. Consequently, realistic exercise edits changed its terminal forecast
by only approximately $0.003$~mg/dL. The completed experiment therefore used
detector-derived exercise pseudo-labels as retrospective supervision and
trained the scenario decoder to use factual future exercise information.

This stage is distinct from both preceding analyses. Detector construction asks
whether wearable signals locate ambulatory exercise-like episodes, while the
descriptive kernels ask how observed glucose behaves around those episodes.
The present factual-future evaluation asks whether revealing the observed
future exercise path improves prediction, and the planned-scenario evaluation
measures the model's sensitivity to supplied exercise profiles. Neither
analysis identifies the glucose trajectory that would have occurred under a
randomly assigned intervention.

\subsubsection{Exercise scenario representation}

The original five-minute panel was not modified. A copied panel retained all
original HR, steps, timestamps, and other columns, and added nine future
scenario channels: \texttt{exercise\_detected\_active}, the light, moderate,
and vigorous \texttt{exercise\_strain} indicators, the slow-walk, walk,
brisk-walk, and run-like \texttt{exercise\_cadence} indicators, and
\texttt{exercise\_duration\_minutes}. The observed history additionally
contained \texttt{recent\_exercise\_30min},
\texttt{recent\_exercise\_60min}, and
\texttt{recent\_exercise\_120min}. HR and steps remained the existing future
physiological scenario variables. Sub-threshold indicators, confidence,
recovery, elapsed-time, and remaining-time fields were retained for audit or
excluded from model input according to their provenance.

Alignment used participant identity and timezone-aware local timestamps. Model
and detector segment identifiers were not assumed to be numerically identical.
The audit retained 1{,}770 panel-aligned episodes: 1{,}204 in training, 282 in
validation, and 284 in test. Exercise-active rows represented only $0.744\%$
of the 4{,}532{,}158-row copied panel, motivating exercise-anchor weighting and
sampling during training.

Empirical planned-exercise profiles were estimated from training participants
only. For strain class $c$ and horizon $h$, the participant-balanced profile was
\begin{equation}
  \overline{x}^{(c)}_h
  =
  \frac{1}{N_c}
  \sum_{p=1}^{N_c}
  \left[
    \frac{1}{n_{p,c}}
    \sum_{i=1}^{n_{p,c}} x^{(p,i,c)}_h
  \right],
  \label{eq:exercise_empirical_signature}
\end{equation}
where $x_h$ is the HR or step trajectory. Episodes were aligned at the
detector-derived onset, averaged within participant, and then averaged across
participants. Twelve five-minute steps represented the one-hour scenario, and
participant-bootstrap intervals preserved participant clustering. Training
support comprised 82 light episodes from 70 participants, 591 moderate episodes
from 280 participants, and 476 vigorous episodes from 297 participants.

Figure~\ref{fig:detector_exercise_signatures} shows that HR clearly separates
the strain classes, whereas step profiles overlap considerably. Strain thus
captures cardiovascular load more clearly than cadence in this library.
Cadence remains an auxiliary descriptor, and run-like denotes high-cadence
wrist-wearable ambulation rather than confirmed running.

\begin{figure}[H]
  \centering
  \includegraphics[width=0.92\linewidth]{exercise_scenario/figure_1_empirical_exercise_signatures.pdf}
  \caption{%
    \textbf{Training-only empirical exercise signatures.}
    The empirical detector-derived exercise library used as planned scenario
    input. Heart-rate trajectories distinguish strain classes more clearly
    than step trajectories. Profiles were estimated from training participants
    only; bands show participant-bootstrap confidence intervals.%
  }
  \label{fig:detector_exercise_signatures}
\end{figure}

\subsubsection{Scenario-aware retraining objective}

The historical state was fixed before the scenario edit,
\begin{equation}
  h_t = E_{\theta}\left(X_{1:t},S\right),
  \label{eq:exercise_history_state}
\end{equation}
and the scenario-aware forecast was decomposed as
\begin{equation}
  \widehat{\mathbf G}_{t+1:t+H}
  = f_{\mathrm{base}}\left(h_t,T\right)
  + \Delta f_{\mathrm{exercise}}
    \left(h_t,T,A\odot M,M\right).
  \label{eq:exercise_scenario_decomposition}
\end{equation}
Here $A$ contains planned future values, $M$ contains their availability
masks, and $A\odot M$ prevents unknown future values from being treated as
observed. The additive exercise branch was initialized at zero.

The mask distinguishes three states. Forecast only uses $M=0$. Planned no
exercise uses $M=1$ with $a^{\mathrm{exercise}}=0$. Planned exercise uses
$M=1$ with $a^{\mathrm{exercise}}=A^{(c)}$ for
$c\in\{\mathrm{light},\mathrm{moderate},\mathrm{vigorous}\}$. Unknown future
is therefore not equivalent to an explicitly supplied plan of no exercise, and
the primary planned comparison is exercise versus planned no exercise.

The imposed Path B kernel was disabled, \texttt{hr\_exercise=false}, and no
fixed negative sign or ordering loss forced vigorous below moderate or light.
The checkpoint audit found all 87 frozen canonical tensors exactly unchanged;
only the expanded scenario and recent-history pathway was learned.

Training combined a forecast-only pass,
\begin{equation}
  \widehat{\mathbf G}^{\mathrm{FO}}
  = f_{\theta}\left(h_t,A=0,M=0\right),
\end{equation}
with a factual-future pass,
\begin{equation}
  \widehat{\mathbf G}^{\mathrm{FF}}
  = f_{\theta}\left(h_t,A=A^{\mathrm{observed}},
      M=M^{\mathrm{observed}}\right).
\end{equation}
The objective was
\begin{equation}
  \mathcal L
  = \mathcal L_{\mathrm{FO}}
  + \lambda_{\mathrm{FF}}\mathcal L_{\mathrm{FF}}
  + \lambda_{\mathrm{smooth}}
      \left\|\Delta^2\Delta\widehat{\mathbf G}\right\|_2^2
  + \lambda_0\left(\Delta\widehat G_1\right)^2.
  \label{eq:exercise_training_objective}
\end{equation}
The forecast-only loss protected ordinary forecasting, the factual-future loss
trained use of observed future exercise, the smoothness term discouraged
oscillatory effects, and the first-step penalty discouraged an instantaneous
jump. Exercise anchors were oversampled or weighted because they were sparse.
The objective did not reward a large scenario-effect magnitude.

Predictive factual-future gain is defined as
\begin{equation}
  G_{\mathrm{FF}}
  = \mathcal L_{\mathrm{FO,exercise}}
  - \mathcal L_{\mathrm{FF,exercise}}.
  \label{eq:exercise_ff_gain}
\end{equation}
A positive value means that exercise information improves prediction on the
same exercise anchors. It is not an intervention-effect estimate.

\subsubsection{Checkpoint selection and channel audit}

The guarded validation criterion selected epoch 0, the first completed
training epoch, from the selected best-model checkpoint. Training continued
through epoch 5 and stopped normally when early-stopping patience reached five;
the stop was not caused by an error. At selection, forecast-only validation MAE
degraded by only $0.0257$~mg/dL relative to the canonical reference, well below
the $0.20$~mg/dL guard, while validation exercise factual-future pinball gain
was $0.0378$. Epoch numbering therefore does not mean that no training occurred.
The new exercise values, availability masks, and recent-history projections
all had nonzero learned weights.

Table~\ref{tab:detector_exercise_support} summarizes the training and evaluation
support. Training and test participants remained disjoint, and the scenario
templates used training participants only.

\begin{table}[H]
  \centering
  \caption{\textbf{Detector-informed training and evaluation support.}}
  \label{tab:detector_exercise_support}
  \small
  \begin{tabular}{lr}
    \toprule
    Quantity & Value \\
    \midrule
    Selected trained epoch & 0 \\
    Eligible held-out exercise episodes & 281 \\
    Held-out exercise participants & 110 \\
    Forecast horizons & 12 \\
    Participant-bootstrap resamples & 500 \\
    Frozen canonical tensors, exactly unchanged & 87 \\
    Full-test participants & 221 \\
    Full-test forecast rows & 1{,}861{,}356 \\
    \bottomrule
  \end{tabular}
\end{table}

\subsubsection{Held-out factual-future forecasting}

The fixed held-out evaluation contained 281 eligible light, moderate, or
vigorous episodes from 110 unseen test participants. Three sub-threshold
episodes were excluded, all remaining anchors matched both checkpoints, and
each episode contributed 12 five-minute horizons. Five hundred bootstrap
resamples used the participant as the sampling unit, retaining all episodes
from a sampled participant together.

As shown in Table~\ref{tab:exercise_factual_future}, revealing the factual
exercise path reduced pinball loss from 4.134 to 3.880, MAE from 12.34 to
11.66~mg/dL, RMSE from 17.63 to 17.01~mg/dL, and 60-minute MAE from 16.22 to
15.18~mg/dL. The predictive factual-future pinball gain was $0.254$ with 95\%
participant-bootstrap CI $[0.174,0.325]$. Forecast-only and factual predictions
differed by 2.09~mg/dL on average. This demonstrates an active, predictively
useful scenario pathway, not the glucose response to an assigned intervention.
Exercise-anchor errors are not compared with full-stream errors because the
evaluation cohorts differ.

\begin{table}[H]
  \centering
  \caption{%
    \textbf{Held-out factual-future forecasting.}
    Intervals are participant-bootstrap 95\% CIs. The gain row is forecast
    only minus factual future, so positive values indicate improvement.%
  }
  \label{tab:exercise_factual_future}
  \scriptsize
  \resizebox{\linewidth}{!}{%
  \begin{tabular}{lcccc}
    \toprule
    Mode & Pinball & MAE, mg/dL & RMSE, mg/dL & 60-min MAE, mg/dL \\
    \midrule
    Forecast only & 4.134 $[3.667,4.622]$ & 12.34 $[11.09,13.78]$ & 17.63 $[15.44,20.42]$ & 16.22 $[14.25,18.39]$ \\
    Factual future exercise & 3.880 $[3.466,4.409]$ & 11.66 $[10.42,12.81]$ & 17.01 $[15.00,19.30]$ & 15.18 $[13.21,17.41]$ \\
    Gain & 0.254 $[0.174,0.325]$ & 0.673 $[0.453,0.853]$ & 0.623 $[0.330,0.856]$ & 1.033 $[0.667,1.350]$ \\
    \bottomrule
  \end{tabular}}
\end{table}

Figure~\ref{fig:detector_exercise_main_results} places the predictive gain next
to pathway activation, planned trajectories, and dependence on starting
glucose. Detailed component figures are reported in
Appendix~\ref{subsec:exercise_scenario_supplement}.

\begin{figure}[H]
  \centering
  \includegraphics[width=0.98\linewidth]{exercise_scenario/figure_8_exercise_result_summary.pdf}
  \caption{%
    \textbf{Detector-informed exercise pathway: main results.}
    Predictive factual-future gain, canonical versus trained sensitivity,
    planned-exercise trajectories, and dependence on starting glucose are
    shown on the identical held-out anchor cohort.%
  }
  \label{fig:detector_exercise_main_results}
\end{figure}

\subsubsection{Learned planned-exercise responses}

For strain class $c$, the model-implied planned response was
\begin{equation}
  \Delta\widehat{\mathbf G}^{(c)}
  = \widehat{\mathbf G}\left(h_t,A^{(c)},M=1\right)
  - \widehat{\mathbf G}\left(h_t,A^{(0)},M=1\right),
  \label{eq:planned_exercise_effect}
\end{equation}
where $A^{(0)}$ denotes planned no exercise. At 60 minutes, the learned effects
were $-1.345$~mg/dL for light, $-2.579$~mg/dL for moderate, and
$-2.916$~mg/dL for vigorous scenarios. On identical anchors, canonical
effects were only $-0.0036$, $-0.0033$, and $-0.0033$~mg/dL, respectively
(Table~\ref{tab:exercise_planned_effects}). The learned response grew gradually
over the horizon, with moderate and vigorous scenarios larger than light.
These are model sensitivities to supplied profiles, not physiological response
estimates.

\begin{table}[H]
  \centering
  \caption{%
    \textbf{Planned scenario effects at 60 minutes relative to planned no
    exercise.} Intervals use participant bootstrap resampling.%
  }
  \label{tab:exercise_planned_effects}
  \small
  \begin{tabular}{lcccrr}
    \toprule
    Scenario & Trained effect & 95\% CI & Canonical effect & Anchors & Participants \\
    \midrule
    Light & $-1.345$ & $[-1.399,-1.295]$ & $-0.0036$ & 281 & 110 \\
    Moderate & $-2.579$ & $[-2.654,-2.488]$ & $-0.0033$ & 281 & 110 \\
    Vigorous & $-2.916$ & $[-3.024,-2.823]$ & $-0.0033$ & 281 & 110 \\
    \bottomrule
  \end{tabular}
\end{table}

A controlled strain-label sweep held physiology, duration, and timing fixed
and produced the same ordering, with terminal effects of approximately
$-1.44$, $-2.60$, and $-2.84$~mg/dL. This confirms that the decoder responds
to the supplied strain category, but does not validate the category as an
assigned biological dose.

Starting glucose was stratified below 105~mg/dL, from 105 to 180~mg/dL, and
above 180~mg/dL. Model-implied reductions were larger at higher starting
glucose, indicating context-conditioned decoder behavior. This pattern does
not establish greater biological effectiveness at high glucose and may reflect
the state representation, the residual-current target, or the observational
training distribution. The corresponding descriptive heterogeneity is kept in
Section~\ref{sec:exercise_subgroup} rather than repeated here.

\subsubsection{Functional reliance and robustness}

Grouped ablations measured loss of predictive gain or increased factual MAE
after removing coherent input groups. Figure~\ref{fig:exercise_grouped_ablation}
and Table~\ref{tab:exercise_functional_ablation} show that the full block gained
$0.254$, labels-only retained $0.219$, and HR-and-steps-only retained $0.095$.
Removing all detector labels reduced gain to $0.042$, while removing all
exercise variables eliminated it. The decoder therefore relies on the combined
exercise block, with detector-derived labels contributing more than HR and
steps alone in this trained pathway.

\begin{figure}[H]
  \centering
  \includegraphics[width=0.92\linewidth]{exercise_scenario/figure_i1_grouped_ablation.pdf}
  \caption{%
    \textbf{Grouped scenario ablation importance.}
    Importance is the loss of predictive factual-future gain or increase in
    factual MAE after removal of a coherent exercise-input group.%
  }
  \label{fig:exercise_grouped_ablation}
\end{figure}

\begin{table}[H]
  \centering
  \caption{%
    \textbf{Functional ablations on identical held-out anchors.}
    Terminal and mean absolute effects are in mg/dL.%
  }
  \label{tab:exercise_functional_ablation}
  \scriptsize
  \resizebox{\linewidth}{!}{%
  \begin{tabular}{lrrrr}
    \toprule
    Analysis & Factual-future gain & Factual MAE & Terminal effect & Mean absolute effect \\
    \midrule
    Full & 0.254 & 11.662 & $-3.127$ & 2.092 \\
    HR and steps only & 0.095 & 12.091 & $-1.013$ & 0.688 \\
    Labels only & 0.219 & 11.752 & $-2.691$ & 1.751 \\
    No strain & 0.194 & 11.824 & $-2.347$ & 1.568 \\
    No cadence & 0.201 & 11.803 & $-2.334$ & 1.550 \\
    No duration & 0.215 & 11.768 & $-2.527$ & 1.692 \\
    All detector labels removed & 0.042 & 12.230 & $-0.381$ & 0.286 \\
    All exercise inputs removed & 0.000 & 12.335 & 0.000 & 0.000 \\
    Shuffled strain and cadence & 0.253 & 11.662 & $-3.132$ & 2.096 \\
    Three-step circular time shift & 0.259 & 11.648 & $-3.191$ & 2.104 \\
    \bottomrule
  \end{tabular}}
\end{table}

Individual channels contain redundant information, and grouped ablation changes
both values and availability-mask semantics. These are functional reliance
measures, not causal feature effects. In particular, shuffled strain and
cadence labels preserved nearly the full predictive gain, the circular
time-shift control changed performance little, and exact-variable permutation
effects were modest and uncertain. The trained model relies strongly on the
presence, availability, and activation structure of the detector-derived
exercise block. Evidence that exact strain and cadence identity is necessary
for factual prediction is weaker. Detailed component ablation, permutation,
probe, and latent-space results appear in
Appendix~\ref{subsec:exercise_scenario_supplement}.

Robustness was evaluated without changing the primary cohort definition.
Factual-future gain remained positive for all primary episodes, $0.254$
$[0.178,0.327]$ on 281 anchors from 110 participants; high-confidence episodes,
$0.281$ $[0.196,0.369]$ on 216 anchors from 98 participants; and strict-recovery
episodes, $0.245$ $[0.140,0.361]$ on 128 anchors from 79 participants. The
result is therefore not driven only by low-confidence episodes. The overlapping
intervals do not establish superiority of either restricted subset; the saved
robustness figure is shown in Appendix~\ref{fig:exercise_robustness_supp}.

\subsubsection{Full-test-stream generalization}

The final forecast-only comparison used 1{,}861{,}356 rows from 221 held-out
test participants. Table~\ref{tab:exercise_full_stream} and
Figure~\ref{fig:exercise_full_stream} show that full-stream MAE was practically
preserved, changing from 10.122~mg/dL for the canonical checkpoint to
10.144~mg/dL for the detector-trained checkpoint. Forecast calibration was not
preserved: bias, defined as prediction minus observation, shifted from
$-1.044$ to $-2.502$~mg/dL. Zeroing the three recent-exercise history channels
yielded MAE 10.141~mg/dL and bias $-2.494$~mg/dL, so those variables do not
explain the global downward shift.

\begin{table}[H]
  \centering
  \caption{\textbf{Full test-stream generalization in forecast-only mode.}}
  \label{tab:exercise_full_stream}
  \small
  \begin{tabular}{lrrrr}
    \toprule
    Condition & MAE, mg/dL & Bias, mg/dL & Participants & Forecast rows \\
    \midrule
    Canonical & 10.122 & $-1.044$ & 221 & 1{,}861{,}356 \\
    Detector trained & 10.144 & $-2.502$ & 221 & 1{,}861{,}356 \\
    Detector trained, recent history zero & 10.141 & $-2.494$ & 221 & 1{,}861{,}356 \\
    \bottomrule
  \end{tabular}
\end{table}

\begin{figure}[H]
  \centering
  \includegraphics[width=0.96\linewidth]{exercise_scenario/figure_13_test_stream_history_ablation.pdf}
  \caption{%
    \textbf{Final test-stream detector comparison and recent-history ablation.}
    Overall MAE is nearly unchanged, while the detector-trained checkpoint has
    a stronger negative bias that remains after recent-exercise history is
    zeroed.%
  }
  \label{fig:exercise_full_stream}
\end{figure}

Distance from exercise clarifies the specialization. During exercise,
detector-trained MAE improved from 11.793 to 11.511~mg/dL, while canonical
positive bias of $+1.311$~mg/dL became a modest negative bias of
$-0.822$~mg/dL. From zero to two hours after exercise, MAE worsened from
12.756 to 12.991~mg/dL and negative bias strengthened from $-1.345$ to
$-4.367$~mg/dL. Zeroing recent history improved the latter MAE to
12.788~mg/dL and bias to $-3.070$~mg/dL, removing part but not all of the
degradation. More than two hours from exercise, MAE was nearly identical
(10.066 versus 10.089~mg/dL), yet the negative bias shift remained
($-1.067$ versus $-2.495$~mg/dL). Exercise specialization therefore improved
the intended context and preserved global absolute error, but did not preserve
calibration. These results do not by themselves establish scenario leakage;
the remaining forecast-only bias requires branch-level decomposition.

\subsubsection{Temporal limitations and future model design}

The controlled duration and timing sweeps in
Figure~\ref{fig:exercise_temporal_sweeps} reveal a horizon-local decoder. A
20-minute supplied scenario returns immediately to zero after 20 minutes, and
a 40-minute scenario returns immediately to zero after 40 minutes. Delaying
the supplied exercise delays the response. The model does not propagate a
post-exercise carryover state and therefore performs horizon-local conditioning
rather than dynamic physiological simulation.

\begin{figure}[H]
  \centering
  \begin{subfigure}[t]{0.48\linewidth}
    \centering
    \includegraphics[width=\linewidth]{exercise_scenario/figure_i6_duration_sensitivity.pdf}
    \caption{Duration sweep.}
  \end{subfigure}\hfill
  \begin{subfigure}[t]{0.48\linewidth}
    \centering
    \includegraphics[width=\linewidth]{exercise_scenario/figure_i7_timing_sensitivity.pdf}
    \caption{Timing sweep.}
  \end{subfigure}
  \caption{%
    \textbf{Temporal scenario sensitivity.}
    Effects occur only at decoder horizons where exercise is actively supplied;
    no gradual post-exercise decay is propagated.%
  }
  \label{fig:exercise_temporal_sweeps}
\end{figure}

A future decoder could introduce a dynamic response state,
\begin{align}
  z_{t+h} &= \rho z_{t+h-1} + \beta a_{t+h}, \qquad 0<\rho<1, \\
  \Delta\widehat G_{t+h} &= g_{\phi}\left(h_t,z_{t+h}\right),
  \label{eq:exercise_dynamic_state}
\end{align}
allowing exercise influence to decay after the active interval rather than
ending abruptly.

Several strengths support the proof of concept. Detector outputs were
operationalized as explicit model variables, participant splits remained
disjoint, empirical templates used training participants only, the canonical
near-zero pathway became responsive, real exercise paths improved held-out
prediction, planned strain classes became differentiated, exercise-period MAE
improved, full-test MAE was essentially unchanged, robustness subsets retained
positive gain, and frozen canonical tensors remained unchanged.

The limitations remain consequential. Pseudo-labels are retrospective and
exercise was not randomized. The empirical profiles are templates rather than
assigned interventions, exact strain and cadence identity is less important
than the coherent block and mask structure, and light support is much smaller
than moderate or vigorous support. Brisk-walk and run-like support is too sparse
for separate claims. The ambulatory detector does not capture stationary
cycling, resistance exercise, or isometric activity. The trained checkpoint
has a stronger negative full-stream bias, the current decoder has no learned
carryover, and the one-hour horizon cannot represent delayed responses.
Prospectively selected examples in Appendix~\ref{fig:exercise_examples_supp}
illustrate both useful responses and remaining unobserved drivers, but do not
replace cohort-level evidence.

The first design priority is structural forecast-only invariance,
\begin{equation}
  M=0 \Longrightarrow \Delta f_{\mathrm{exercise}}=0.
  \label{eq:exercise_forecast_invariance}
\end{equation}
The second is a separate audit of the base branch and zero-mask scenario branch
to locate the downward bias. The third is the dynamic exercise-response and
recovery state above, followed by two- or three-hour evaluation and held-out
calibration of the detector-informed checkpoint. A mask-versus-content control
should then compare real values with real masks, zero values with real masks,
exercise-active only, shuffled values with real masks, and zero values with
zero masks. Finally, validation requires prospective records of exercise type,
duration, meals, and treatment timing.
